\begin{titlepage}
\centering
{\large ĐẠI HỌC QUỐC GIA TP. HỒ CHÍ MINH\\ TRƯỜNG ĐẠI HỌC BÁCH KHOA\\ \textbf{KHOA KHOA HỌC VÀ KỸ THUẬT MÁY TÍNH}\par}
\vspace{1cm}
\includegraphics[width=3.5cm]{hcmut-logo}\par
\vspace{1cm}
{\LARGE\textbf{BÁO CÁO BÀI TẬP LỚN}\par}
\vspace{0.3cm}
{\large MÔN HỌC: INTERNET OF THINGS (IoT)\par}
\vspace{1cm}
{\large ĐỀ TÀI\par}
{\Large\textbf{\topicname}\par}
\vspace{1cm}
Giảng viên hướng dẫn: TS. Lê Trọng Nhân\par
Lớp: \classcode\ -- Nhóm: \groupcode\par
\vspace{0.5cm}
\begin{tabular}{|c|l|l|l|}
\hline STT & Họ và tên & MSSV & Email \\ \hline
1 & \todo{Họ và tên} & \todo{MSSV} & \todo{email}@hcmut.edu.vn \\ \hline
2 & \todo{Họ và tên} & \todo{MSSV} & \todo{email}@hcmut.edu.vn \\ \hline
3 & \todo{Họ và tên} & \todo{MSSV} & \todo{email}@hcmut.edu.vn \\ \hline
\end{tabular}
\vfill
TP. Hồ Chí Minh, tháng \todo{…} năm \todo{20…}
\end{titlepage}

\section*{THÔNG TIN NỘP BÀI}
\begin{tabular}{|l|l|}
\hline Hạng mục & Đường dẫn / Thông tin \\ \hline
Kho mã nguồn GitHub & \url{https://github.com/khoaphunsonac/Assignment_IOT_HK261} \\ \hline
Video demo & \todo{https://…} \\ \hline
Dashboard CoreIOT (public link, nếu có) & \todo{https://…} \\ \hline
Ngày nộp & \todo{dd/mm/yyyy} \\ \hline
\end{tabular}

\section*{BẢNG PHÂN CÔNG CÔNG VIỆC}
\begin{tabularx}{\linewidth}{|c|l|l|X|c|}
\hline STT & Họ và tên & MSSV & Công việc phụ trách & Đóng góp \\ \hline
1 & \todo{Họ và tên} & \todo{MSSV} & \todo{Công việc, mã yếu tố} & \todo{…\%} \\ \hline
2 & \todo{Họ và tên} & \todo{MSSV} & \todo{Công việc, mã yếu tố} & \todo{…\%} \\ \hline
3 & \todo{Họ và tên} & \todo{MSSV} & \todo{Công việc, mã yếu tố} & \todo{…\%} \\ \hline
\end{tabularx}

\clearpage
\section*{BẢNG TỰ ĐÁNH GIÁ CÁC YẾU TỐ CÔNG NGHỆ}
{\small
\textbf{Bảng A. Các yếu tố công nghệ (tối đa 8 điểm)}
\begin{longtable}{|l|p{5.6cm}|c|l|l|l|}
\hline Mã & Yếu tố & Điểm & Hiện thực & Mục & Minh chứng \\ \hline
\endhead
\multicolumn{6}{|l|}{\textbf{Nhóm 1 – Cảm biến}} \\ \hline
T01 & Cảm biến sử dụng giao tiếp I2C trên Tasmota & 0.5 & \todo{Có/Không} & \todo{x.y} & \todo{link} \\ \hline
T02 & Cảm biến sử dụng giao tiếp Modbus RS485 & 1.0 & \todo{Có/Không} & \todo{x.y} & \todo{link} \\ \hline
T03 & Cảm biến sử dụng kết nối BLE với thiết bị khác & 1.5 & \todo{Có/Không} & \todo{x.y} & \todo{link} \\ \hline
\multicolumn{6}{|l|}{\textbf{Nhóm 2 – Trí tuệ nhân tạo trên thiết bị (TinyML)}} \\ \hline
T04 & Thông tin được nội suy từ AI (TinyML chạy trên thiết bị) & 1.0 & \todo{Có/Không} & \todo{x.y} & \todo{link} \\ \hline
T05 & Công cụ thu thập dữ liệu, huấn luyện và export ra model TinyML & 1.0 & \todo{Có/Không} & \todo{x.y} & \todo{link} \\ \hline
\multicolumn{6}{|l|}{\textbf{Nhóm 3 – Giao tiếp và giao diện cục bộ}} \\ \hline
T06 & Giao tiếp không dây sử dụng ESP-NOW & 1.5 & \todo{Có/Không} & \todo{x.y} & \todo{link} \\ \hline
T07 & Webserver trên thiết bị (tự hiện thực – hiển thị, điều khiển local) & 1.0 & \todo{Có/Không} & \todo{x.y} & \todo{link} \\ \hline
\multicolumn{6}{|l|}{\textbf{Nhóm 4 – IoT Gateway cục bộ}} \\ \hline
T08 & IoT Gateway hiện thực trên PC (hoặc Raspberry Pi): tổng hợp dữ liệu, gửi lên CoreIOT, nhận lệnh điều khiển từ CoreIOT và phân phối cho các thiết bị trong mạng local & 1.0 & \todo{Có/Không} & \todo{x.y} & \todo{link} \\ \hline
T09 & Chức năng AI chạy trên IoT Gateway local để phân tích dữ liệu, phát hiện (detect) và gửi kết quả lên server & 1.0 & \todo{Có/Không} & \todo{x.y} & \todo{link} \\ \hline
T10 & Giao diện giám sát – điều khiển trên IoT Gateway local & 1.0 & \todo{Có/Không} & \todo{x.y} & \todo{link} \\ \hline
\multicolumn{6}{|l|}{\textbf{Nhóm 5 – Nền tảng CoreIOT}} \\ \hline
T11 & Sử dụng template để phát triển ứng dụng, thiết kế dashboard, giám sát và điều khiển từ xa & 1.0 & \todo{Có/Không} & \todo{x.y} & \todo{link} \\ \hline
T12 & Sử dụng RuleChain để phát triển tính năng thông minh & 1.0 & \todo{Có/Không} & \todo{x.y} & \todo{link} \\ \hline
T13 & Gửi cảnh báo bằng notification khi dữ liệu vượt ngưỡng & 0.5 & \todo{Có/Không} & \todo{x.y} & \todo{link} \\ \hline
T14 & Cập nhật firmware OTA cho thiết bị & 2.0 & \todo{Có/Không} & \todo{x.y} & \todo{link} \\ \hline
T15 & IoT Gateway trên CoreIOT đồng bộ với IoT Gateway local & 2.0 & \todo{Có/Không} & \todo{x.y} & \todo{link} \\ \hline
\end{longtable}
\textbf{Bảng B. Các tính năng nâng cao (tối đa 1 điểm)}
\begin{longtable}{|l|p{5.6cm}|c|l|l|l|}
\hline Mã & Yếu tố & Điểm & Hiện thực & Mục & Minh chứng \\ \hline
\endhead
N01 & Trợ lý ảo – LLM & 1.0 & \todo{Có/Không} & \todo{x.y} & \todo{link} \\ \hline
N02 & Giọng nói – Âm thanh & 1.0 & \todo{Có/Không} & \todo{x.y} & \todo{link} \\ \hline
N03 & Màn hình hiển thị dựa trên LVGL & 1.0 & \todo{Có/Không} & \todo{x.y} & \todo{link} \\ \hline
N04 & AI chạy bằng GPU của thiết bị & 1.0 & \todo{Có/Không} & \todo{x.y} & \todo{link} \\ \hline
N05 & Ứng dụng mobile được thiết kế riêng & 1.0 & \todo{Có/Không} & \todo{x.y} & \todo{link} \\ \hline
N06 & Giao diện dashboard thiết kế riêng, có khả năng phát hiện (detect) và ghi nhận (record) các sự kiện & 1.0 & \todo{Có/Không} & \todo{x.y} & \todo{link} \\ \hline
N07 & Tự triển khai (deploy) ThingsBoard trên máy local & 1.0 & \todo{Có/Không} & \todo{x.y} & \todo{link} \\ \hline
\end{longtable}
}
\textbf{Tổng điểm tự đánh giá:} Công nghệ: \todo{…}/8 -- Nâng cao: \todo{…}/1

\clearpage
\tableofcontents

\section*{DANH MỤC TỪ VIẾT TẮT}
\begin{tabular}{|l|l|}
\hline Từ viết tắt & Ý nghĩa \\ \hline
BLE & Bluetooth Low Energy \\ \hline
OTA & Over-The-Air (cập nhật firmware từ xa) \\ \hline
TinyML & Tiny Machine Learning \\ \hline
\end{tabular}
